\subsection{Hyperjump за \( \mathcal{O}(N) \)}

\subsubsection{Описание решения}

Переформулировав условие задачи, для данного массива нужно найти отрезок с максимальной суммой его элементов.

Воспользуемся жадной стратегией: будем набирать элементы в отрезок до тех пор, пока их сумма неотрицательная. Иначе выгоднее начать новый отрезок: если продолжить набирать элементы в отрезок с отрицательной суммой, то мы будем в проигрыше на эту сумму.

Остаётся на каждой итерации фиксировать текущую сумму и обновлять глобальный максимум при необходимости.

\subsubsection{Асимптотическая сложность}

\textit{Временная сложность решения} --- линейная \( \mathcal{O}(N) \). Для подсчёта требуемой суммы решение считывает весь массив поэлементно и выполняет для каждого из них константное число операций.

\textit{Пространственная сложность решения} --- константная \( \mathcal{O}(1) \). Алгоритм использует для вычисления результата фиксированное число переменных, исходный массив в памяти не сохраняется.

\subsubsection{Листинг кода}

\begin{lstlisting}
#include <iostream>

int main() {
  int N;
  std::cin >> N;
  long long ans = 0, curr = 0;
  for (int i = 0; i < N; i++) {
    int item;
    std::cin >> item;
    curr += item;
    if (curr < 0) {
      curr = 0;
    }
    if (curr > ans) {
      ans = curr;
    }
  }
  std::cout << ans;
  return 0;
}
\end{lstlisting}